\documentclass[UTF8,a4paper,11pt]{ctexart}
\usepackage[slantfont,boldfont]{xeCJK}
\usepackage{fontspec}


\usepackage{mathtools}
\usepackage{amsmath}
\usepackage{amsfonts}
\usepackage{amssymb}
\usepackage{amsthm}
\usepackage{indentfirst}
\usepackage{graphicx}
\usepackage{geometry}
\usepackage{fancyhdr}
\usepackage{titlesec}
\usepackage{setspace}
\usepackage{enumitem}
\usepackage{bm}
\usepackage{hyperref}

\setlength{\parindent}{2em}
\linespread{1.2}
\geometry{left=2.5cm,right=2.5cm,top=2.5cm,bottom=2.5cm}

\everymath{\displaystyle}

\newcommand*{\QED}{\hfill\ensuremath{\square}}

\setlength{\headheight}{13.6pt}
\pagestyle{fancy}
\fancyhf{}
\fancyfoot[C]{\thepage}
\fancyhead[L]{数值分析\quad 第一章理论作业\quad 学号：241098117，姓名：张城宗}
\fancyhead[R]{习题一}

\begin{document}

\begin{center}
  {\LARGE\bfseries 数值分析\quad 第一章理论作业}\\[6pt]
  {\large 习题一\quad 第 1、3、4、5、7、9、10 题}\\[4pt]
  {\normalsize 学号：241098117\quad 姓名：张城宗}
\end{center}

\bigskip

%=====================================================
\noindent\textbf{第 1 题.}\quad 应用梯形公式
\[
  T = \frac{b-a}{2}\bigl(f(a)+f(b)\bigr)
\]
计算积分
\[
  I = \int_0^1 e^{-x}\,\mathrm{d}x
\]
的近似值，在整个计算过程中按四舍五入规则取五位小数．计算中产生的误差的主要原因是离散或是舍入？为什么？

\medskip
\noindent\textbf{解.}\quad 取 $a=0,\ b=1,\ f(x)=e^{-x}$，有
\[
  f(0) = 1.00000,\qquad f(1) = e^{-1} \approx 0.36788.
\]
代入梯形公式：
\[
  T = \frac{1}{2}(1.00000 + 0.36788) = 0.68394.
\]
精确值为 $I = 1 - e^{-1} \approx 0.63212$，故离散误差
\[
  |T - I| \approx 0.05182 \sim O(10^{-2}).
\]
计算中仅 $f(1)$ 的舍入引入不超过 $\frac{1}{2}\times 10^{-5}$ 的误差，乘以 $\frac{1}{2}$ 后量级仍为
$O(10^{-5})$，远小于离散误差．

故\textbf{主要误差来源是离散误差}．原因在于：梯形公式以端点连线代替曲线，而
$e^{-x}$ 在 $[0,1]$ 上是凸函数，弦始终在曲线上方，造成系统性的方法偏差；五位小数的精度对于抑制舍入误差已绰绰有余．

\bigskip\bigskip

%=====================================================
\noindent\textbf{第 3 题.}\quad 证明：若近似数
$\bar{a} = \pm a_0 a_1\cdots a_m.a_{m+1}\cdots a_n$
的相对误差有估计式
\[
  \left|\frac{a-\bar{a}}{\bar{a}}\right| \leqslant \frac{1}{2(a_s+1)}\times 10^{-(n-s)},
\]
其中 $a_s\neq 0$ 是 $\bar{a}$ 的第一位有效数字，则 $\bar{a}$ 至少具有 $n+1-s$ 位有效数字．

\medskip
\noindent\textbf{证明.}\quad
因 $a_0=a_1=\cdots=a_{s-1}=0$，$a_s\neq 0$，由 $\bar{a}$ 的十进制展开知
\[
  |\bar{a}| < (a_s+1)\times 10^{m-s}.
\]
于是
\begin{align*}
  |a-\bar{a}|
  &= |\bar{a}|\cdot\left|\frac{a-\bar{a}}{\bar{a}}\right|
   \leqslant |\bar{a}|\cdot\frac{1}{2(a_s+1)}\times 10^{-(n-s)}\\
  &< (a_s+1)\times 10^{m-s}\cdot\frac{1}{2(a_s+1)}\times 10^{-(n-s)}
   = \frac{1}{2}\times 10^{-(n-m)}.
\end{align*}
由式 (1.2.2) 的定义，$|a-\bar{a}|<\dfrac{1}{2}\times 10^{-(n-m)}$ 且首位非零数字为 $a_s$，故 $\bar{a}$ 至少具有 $n+1-s$ 位有效数字．\QED

\bigskip\bigskip

%=====================================================
\noindent\textbf{第 4 题.}\quad 设 $\bar{x}=23.3123,\ \bar{y}=23.3122$，且
$|e_{\bar{x}}|\leqslant\dfrac{1}{2}\times 10^{-4},\ |e_{\bar{y}}|\leqslant\dfrac{1}{2}\times 10^{-4}$．
问差 $\bar{x}-\bar{y}$ 最多有几位有效数字？

\medskip
\noindent\textbf{解.}\quad
差的近似值 $\bar{x}-\bar{y}=0.0001=10^{-4}$，绝对误差界为
\[
  |e_{\bar{x}-\bar{y}}|
  = \bigl|(x-\bar{x})-(y-\bar{y})\bigr|
  \leqslant |e_{\bar{x}}|+|e_{\bar{y}}|
  \leqslant \frac{1}{2}\times 10^{-4}+\frac{1}{2}\times 10^{-4} = 10^{-4}.
\]
$\bar{x}-\bar{y}$ 的首位有效数字 $1$ 位于 $10^{-4}$ 处，保证 $1$ 位有效数字须 $|e|\leqslant\frac{1}{2}\times 10^{-4}$，
而上界 $10^{-4}>\frac{1}{2}\times 10^{-4}$，条件不满足．故\textbf{差不能保证有任何有效数字}，这是典型的相减相消现象：两数仅末位不同，相减后有效位全部丢失．

\bigskip\bigskip

%=====================================================
\noindent\textbf{第 5 题.}\quad 当 $x(>0)$ 很大时，如何计算
(1)~$\arctan(x+1)-\arctan x$；\quad
(2)~$\ln(x-\sqrt{x^2-1})$；\quad
(3)~$\dfrac{\tan x}{x-\sqrt{x^2-1}}$，
可使其误差较小？

\medskip
\noindent\textbf{解.}\quad 当 $x$ 很大时，以上三式的原始形式均含有两个相近量相减，会发生严重的相减消去，
应通过代数恒等变形消去病态的减法运算．

\begin{enumerate}[label=(\arabic*),leftmargin=*,itemsep=6pt]
  \item 利用反正切差化积公式
        \[
          \arctan(x+1)-\arctan x
          = \arctan\frac{(x+1)-x}{1+x(x+1)}
          = \arctan\frac{1}{1+x(x+1)}.
        \]
        当 $x$ 很大时，$1+x(x+1)\approx x^2$，右端自变量接近零，$\arctan$ 可精确计算，
        避免了原式中两个均趋于 $\dfrac{\pi}{2}$ 的大量相减．

  \item 分子有理化：将 $x-\sqrt{x^2-1}$ 乘除以其共轭式，
        \[
          x-\sqrt{x^2-1}
          = \frac{(x-\sqrt{x^2-1})(x+\sqrt{x^2-1})}{x+\sqrt{x^2-1}}
          = \frac{x^2-(x^2-1)}{x+\sqrt{x^2-1}}
          = \frac{1}{x+\sqrt{x^2-1}},
        \]
        故
        \[
          \ln(x-\sqrt{x^2-1}) = \ln\frac{1}{x+\sqrt{x^2-1}} = -\ln(x+\sqrt{x^2-1}).
        \]
        右端 $x+\sqrt{x^2-1}$ 为两个大量之和，不发生消去，可精确计算．

  \item 同样对分母作有理化：
        \begin{align*}
          \frac{\tan x}{x-\sqrt{x^2-1}}
          &= \frac{\tan x \cdot (x+\sqrt{x^2-1})}{(x-\sqrt{x^2-1})(x+\sqrt{x^2-1})} \\
          &= \frac{\tan x \cdot (x+\sqrt{x^2-1})}{x^2-(x^2-1)}
           = \tan x\cdot(x+\sqrt{x^2-1}).
        \end{align*}
        右端分母消去，改用 $x+\sqrt{x^2-1}$ 的和式，避免了小量除法引起的误差放大．
\end{enumerate}

\bigskip\bigskip

%=====================================================
\noindent\textbf{第 7 题.}\quad 在 $p=10,\,t=4,\,L=U=5$ 的采用舍入近似的假想计算机上，将下列实数写成相应的规格化浮点数：
(1)~$-315.37$；\quad (2)~$0.0074$；\quad (3)~$25\dfrac{2}{3}$；\quad (4)~$52\dfrac{1}{3}$．

\medskip
\noindent\textbf{解.}\quad 规格化浮点形式为 $\pm 0.d_1 d_2 d_3 d_4\times 10^J$（$d_1\neq 0$，$-5\leqslant J\leqslant 5$），
尾数保留四位，按舍入取近似．

\begin{enumerate}[label=(\arabic*),leftmargin=*,itemsep=4pt]
  \item $-315.37 = -0.31537\times 10^3$，第五位尾数为 $7\geqslant 5$，末位进一：
        \[
          fl(-315.37) = -0.3154\times 10^3.
        \]

  \item $0.0074 = 0.7400\times 10^{-2}$，尾数恰好四位，精确表示：
        \[
          fl(0.0074) = 0.7400\times 10^{-2}.
        \]

  \item $25\dfrac{2}{3} = 25.6\overline{6} = 0.256\overline{6}\times 10^2$，第五位尾数为 $6\geqslant 5$，末位进一：
        \[
          fl\!\left(25\tfrac{2}{3}\right) = 0.2567\times 10^2.
        \]

  \item $52\dfrac{1}{3} = 52.\overline{3} = 0.52\overline{3}\times 10^2$，第五位尾数为 $3<5$，截断：
        \[
          fl\!\left(52\tfrac{1}{3}\right) = 0.5233\times 10^2.
        \]
\end{enumerate}

\bigskip\bigskip

%=====================================================
\noindent\textbf{第 9 题.}\quad 假设在 $p=10,\,t=4,\,L=U=5$ 的采用舍入近似的假想计算机上作浮点运算．设
$x = 0.6436\times 10^{-4},\ y = 0.8321\times 10^{-3}$，求 $x\oplus y$ 和 $x\ominus y$．

\medskip
\noindent\textbf{解.}\quad 对阶，将 $x$ 调整至指数 $-3$：$x = 0.06436\times 10^{-3}$．

加法：$0.06436 + 0.8321 = 0.89646$，第五位为 $6\geqslant 5$，末位进一，故
\[
  x\oplus y = 0.8965\times 10^{-3}.
\]

减法：$0.06436 - 0.8321 = -0.76774$，第五位为 $4<5$，截断，故
\[
  x\ominus y = -0.7677\times 10^{-3}.
\]

\bigskip\bigskip

%=====================================================
\noindent\textbf{第 10 题.}\quad 设字长 $t=8$，以及
\[
  x = 0.32282167\times 10^{-4},\quad
  y = 0.45367549\times 10^{2},\quad
  z = -0.45368822\times 10^{2}.
\]
求 $fl(fl(x+y)+z)$，$fl(x+fl(y+z))$ 和 $x+y+z$．

\medskip
\noindent\textbf{解.}

\smallskip
\noindent\textbf{(1)~$fl(fl(x+y)+z)$.}\quad
先求 $fl(x+y)$：$x$ 的指数为 $-4$，$y$ 的指数为 $2$，相差六位，对阶后
\[
  x = 0.0000000032282167\times 10^2,
\]
求和 $0.4536754932\cdots\times 10^2$，第九位为 $3<5$，截断，得
\[
  fl(x+y) = 0.45367549\times 10^2 = y.
\]
$x$ 因量级过小被完全吸收，信息丢失．再与 $z$ 求和：
\[
  y + z = (0.45367549 - 0.45368822)\times 10^2 = -0.00001273\times 10^2,
\]
规格化后恰为八位有效数字，无需再舍入：
\[
  fl\bigl(fl(x+y)+z\bigr) = -0.12730000\times 10^{-2}.
\]

\smallskip
\noindent\textbf{(2)~$fl(x+fl(y+z))$.}\quad
先求 $fl(y+z)$：$y$ 与 $z$ 同阶，
\[
  y + z = -0.00001273\times 10^2 = -0.12730000\times 10^{-2},
\]
恰为八位，无舍入，故 $fl(y+z)=-0.12730000\times 10^{-2}$．再加 $x$，对阶：
\[
  x = 0.0032282167\times 10^{-2},
\]
求和 $0.0032282167+(-0.12730000) = -0.1240717833$，第九位为 $3<5$，截断：
\[
  fl\bigl(x+fl(y+z)\bigr) = -0.12407178\times 10^{-2}.
\]

\smallskip
\noindent\textbf{(3)~精确值.}\quad
\[
  x+y+z
  = 0.0032282167\times 10^{-2} + (-0.12730000\times 10^{-2})
  = -0.1240717833\times 10^{-2}.
\]

\smallskip
\noindent\textbf{小结.}\quad 先算 $fl(x+y)$ 时，$x\ll y$，对阶使 $x$ 的八位有效数字全部超出
尾数范围而被截断，最终误差约 $3.2\times 10^{-4}$（相对误差约 $26\%$）；
先算 $fl(y+z)$ 时 $y,z$ 量级相当，消去大数后 $x$ 可正常参与运算，
结果仅末位有舍入误差，与精确值相差 $5\times 10^{-10}$．
两种顺序精度相差悬殊，说明运算次序对舍入误差的积累有实质影响．

\end{document}
